\documentclass[10pt]{article}
\usepackage[a4paper,margin=2cm,top=2.4cm,bottom=2cm]{geometry}
\usepackage{amsmath,amssymb}
\usepackage{fancyhdr}
\usepackage{enumitem}
\usepackage{titlesec}
\usepackage[hidelinks]{hyperref}

% Edit these two macros to match the program's header.
\newcommand{\schoolname}{Efi Arazi School of Computer Science \quad Machine Learning and Data Science M.Sc. Program}
\newcommand{\proposalname}{MLDS Project Proposal}

\pagestyle{fancy}
\fancyhf{}
\renewcommand{\headrulewidth}{0.4pt}
\fancyhead[L]{\footnotesize\schoolname}
\fancyfoot[R]{\footnotesize\proposalname\quad Page \thepage}

\setlength{\parindent}{0pt}
\setlength{\parskip}{0.4em}
\setlist{itemsep=0.1em,topsep=0.2em,leftmargin=1.3em}
\titleformat{\section}{\normalfont\large\bfseries}{}{0em}{}
\titlespacing*{\section}{0pt}{0.8em}{0.3em}

\begin{document}

{\Large\bfseries \proposalname\par}
\vspace{0.1em}
{\large\bfseries Label-Free Aggregation of Human Annotators with FUSE, and Training on the
Result\par}
\vspace{0.4em}

\textbf{Date:} November 2026 \qquad \textbf{Mentors:} TBD

\section*{Background}

Supervised models are only as good as their labels, and labels from a single expert are
slow, costly, and still imperfect. A common alternative is to collect several cheaper
annotations per item and combine them. Majority vote treats all annotators as equally
reliable; Dawid--Skene \cite{ds} learns each annotator's reliability by EM but can be
unstable when each item has few annotations. FUSE (Fully Unsupervised Score Ensembling)
\cite{fuse} takes a different route: it treats each annotator as a noisy binary classifier,
estimates its sensitivity and specificity by the method of moments of Jaffe et
al.~\cite{jaffe} from the agreement structure among annotators alone, and outputs a
posterior probability per item, with no labels at all. FUSE was designed and evaluated for
ranking language-model answers using LLM judges as the ``annotators.'' Whether it helps
when the annotators are people, and whether the labels it produces train a better model
than one expert's labels, is open. This project asks both.

\textbf{Preliminaries.} Item $i$ has an unobserved label $y_i\in\{1,\dots,K\}$; annotator
$j$ labels a subset of items, giving $A_{ij}$ or a missing entry. In the binary case, annotator
$j$ has sensitivity $\psi_j$ and specificity $\eta_j$; under conditional independence given
$y$ the off-diagonal covariance of $A$ is rank one with entries $\lambda_j\lambda_k$,
$\lambda_j\propto\psi_j+\eta_j-1$, so each $\lambda_j$ is identifiable from any three
annotators. Multi-class tasks are handled one-vs-rest. FUSE's posterior $p_i$ is used as a
hard label ($\arg\max$) or as a soft training target. A ``single expert'' means one
annotator's label on every training item. Human annotators may violate conditional
independence, since they share guidelines, biases and class priors, so independence is a
risk to be measured, not assumed.

\textbf{Stage 1 (until the midterm presentation): does FUSE improve the aggregation
decision on human labels?}
\begin{enumerate}[label=\alph*)]
\item \textbf{Data.} Four to six public datasets with several human annotations per item
and a curated gold label: binary (textual entailment, Snow et al.~\cite{snow}; relevance
and sentiment sets from the Zheng et al.\ benchmark~\cite{zheng}) and multi-class
(CIFAR-10H~\cite{c10h}; LabelMe~\cite{labelme}). Gold labels are hidden from every method
and used only for scoring.
\item \textbf{Implementation.} Implement FUSE from the paper and extend it to what human
data needs: missing entries (pairwise-complete moments with a minimum-overlap rule),
multi-class, and calibrated posteriors. Unit-test on simulated annotators with known
$\psi_j,\eta_j$.
\item \textbf{Aggregation benchmark} (mandatory). FUSE against majority vote, Dawid--Skene
and Jaffe et al., on accuracy and log loss against gold. Annotations per item $k$ are
subsampled to trace accuracy as a function of annotation budget.
\item \textbf{Stress tests} (secondary). Simulated and subsampled annotator groups with
shared bias, spammers and adversarial minorities, to find where the method degrades, and
whether a label-free check can flag it.
\end{enumerate}

\textbf{Midterm criterion, fixed in advance.} On a dataset, Stage 1 succeeds if FUSE's
accuracy gain over majority vote has a paired-bootstrap 95\% interval above zero at a
stated $k$; on the study, if this holds on most datasets. ``FUSE beats majority vote only
in regime X'' is a valid outcome, and Stage 2 then uses the best aggregator found.

\textbf{Stage 2: does a model trained on FUSE labels beat one trained on a single
expert?} A classifier (small CNN on image sets; pretrained encoder on a text set) is
trained with identical architecture, optimiser and budget on labels from: (1) a single
annotator, in three versions (a random annotator; the best annotator chosen by gold
agreement, an optimistic bound for ``one expert''; the dataset's own expert where one
exists); (2) majority vote; (3) Dawid--Skene; (4) FUSE hard labels; (5) FUSE soft labels;
(6) gold labels as the ceiling. Each is scored on a held-out gold-labelled test split no
aggregator touched, by accuracy, log loss and calibration error over at least five seeds.
Because several crowd labels cost more than one, the comparison is also made at equal
annotation cost over a range of assumed expert-to-crowd price ratios, and the ratio where
the advantage disappears is reported. \emph{Extension (gated on a Stage 2 gain):} feed the
trained model back into FUSE as one more annotator.

\textbf{Evaluation discipline.} \emph{Leakage:} gold enters only in scoring, and
hyperparameters are tuned on the single-annotator arm so FUSE is never the favoured arm.
\emph{Circularity:} aggregated labels train on one item split and are tested on a disjoint
one. \emph{Imperfect gold:} gold labels are noisy, so near-ceiling accuracy is not
interpreted. \emph{Multiple comparisons:} results are paired bootstrap intervals over
items, with every dataset, arm and $k$ reported, not the best cell. \emph{Observational
limits:} datasets differ in annotator recruitment, so differences are described, not
causally attributed.

\section*{Objectives \& Deliverables}

\textbf{Objectives.} (1) Adapt FUSE to human annotators. (2) Show, against gold and under
the criterion above, whether FUSE improves label aggregation over standard methods and at
which annotation budgets; this closes Stage 1. (3) Train models on FUSE labels and compare
them with models trained on one expert, at equal items and at equal cost; this is the
central contribution.

\textbf{Deliverables.} An open-source FUSE library with tests for missing data and
multi-class tasks; a reproducible benchmark (aggregation accuracy by $k$; Stage 2 metrics by
arm with intervals); a midterm presentation evaluating the Stage 1 criterion; a final report
stating, falsifiably, when FUSE labels beat a single expert's, when they do not, and the
cost ratio at which the answer flips.

\section*{Requirements \& Resources}

Advanced statistics, Python and a deep-learning framework; about 1--2 days a week for an
academic year. Public datasets only; Stage 1 runs on a laptop, Stage 2 on a single GPU.
Academic project; any IP will be addressed when such occurs.

\begin{thebibliography}{9}\footnotesize\setlength{\itemsep}{0pt}
\bibitem{fuse} Lee, Ma, Zhao, et al. Fully Unsupervised Score Ensembling (FUSE). arXiv:2604.18547.
\bibitem{jaffe} Jaffe, A., Nadler, B., Kluger, Y. (2015). Estimating the accuracies of multiple classifiers without labeled data. AISTATS.
\bibitem{ds} Dawid, A. P., Skene, A. M. (1979). Maximum likelihood estimation of observer error-rates using the EM algorithm. J. Roy. Stat. Soc. C 28(1), 20--28.
\bibitem{snow} Snow, R., O'Connor, B., Jurafsky, D., Ng, A. Y. (2008). Cheap and fast---but is it good? Evaluating non-expert annotations for natural language tasks. EMNLP.
\bibitem{zheng} Zheng, Y., Li, G., Li, Y., Shan, C., Cheng, R. (2017). Truth inference in crowdsourcing: Is the problem solved? PVLDB 10(5).
\bibitem{c10h} Peterson, J. C., Battleday, R. M., Griffiths, T. L., Russakovsky, O. (2019). Human uncertainty makes classification more robust. ICCV.
\bibitem{labelme} Rodrigues, F., Pereira, F. C. (2018). Deep learning from crowds. AAAI.
\end{thebibliography}

\end{document}
